\documentclass[10pt,a4paper]{article}

% ============================================================
% PAQUETES
% ============================================================
\usepackage[utf8]{inputenc}
\usepackage[T1]{fontenc}
\usepackage{lmodern}
\usepackage[spanish,provide=*]{babel}
\usepackage[a4paper,left=1.7cm,right=1.7cm,top=1.4cm,bottom=1.2cm]{geometry}
\usepackage{xcolor}
\usepackage{titlesec}
\usepackage{enumitem}
\usepackage{hyperref}

\ifdefined\pdfgentounicode
    \pdfgentounicode=1
\fi

% Tipografía sans-serif (más moderna y legible en pantalla)
\renewcommand{\familydefault}{\sfdefault}

% ============================================================
% COLORES (un solo acento)
% ============================================================
\definecolor{primario}{HTML}{0F766E}
\definecolor{secundario}{HTML}{64748B}
\definecolor{texto}{HTML}{1E293B}
\color{texto}

\hypersetup{colorlinks=true, linkcolor=primario, urlcolor=primario}

% ============================================================
% SECCIONES: título pequeño en mayúsculas + línea fina
% ============================================================
\titleformat{\section}
    {\normalfont\small\bfseries\color{primario}}
    {}{0em}{}
    [\vspace{-3pt}{\color{primario}\titlerule[0.6pt]}]
\titlespacing*{\section}{0pt}{11pt}{5pt}

% ============================================================
% COMANDOS
% ============================================================
% \entrada{Título}{Fecha}{Subtítulo}
\newcommand{\entrada}[3]{%
    \noindent\textbf{#1}\hfill{\small\color{secundario}#2}\\[-1pt]
    {\small\color{secundario}#3}\par\vspace{2pt}%
}
% \fila{Etiqueta}{Contenido}  (para habilidades)
\newcommand{\fila}[2]{%
    \noindent\textbf{\small\color{primario}#1}\hspace{6pt}{\small #2}\par\vspace{2.5pt}%
}

\setlist[itemize]{leftmargin=11pt, topsep=1pt, itemsep=1.5pt, parsep=0pt,
    label=\textcolor{primario}{\scriptsize$\bullet$}}
\setlength{\parskip}{0pt}
\setlength{\parindent}{0pt}
\pagestyle{empty}

\begin{document}

% ============================================================
% ENCABEZADO
% ============================================================
\begin{center}
    {\fontsize{21}{24}\selectfont\bfseries\color{primario} Daniel Fernando Aranda Contreras}\\[4pt]
    {\small\color{secundario} Estudiante de Ingeniería Eléctrica \quad$\cdot$\quad Universidad Industrial de Santander}\\[4pt]
    {\small
        Bucaramanga, Colombia \quad$\cdot$\quad +57 3134996076 \quad$\cdot$\quad
        \href{mailto:dfac24138@gmail.com}{dfac24138@gmail.com} \quad$\cdot$\quad
        \href{https://github.com/ferniad}{github.com/ferniad}
    }
\end{center}

% ============================================================
% PERFIL
% ============================================================
\section*{PERFIL}
Estudiante de 9.\textsuperscript{o} semestre de Ingeniería Eléctrica con base en sistemas de potencia, energía solar fotovoltaica y automatización industrial. Estimo \textit{Hosting Capacity} con \textbf{OpenDSS y Python}. Busco una pasantía en optimización de redes, integración solar o analítica energética.

% ============================================================
% FORMACIÓN
% ============================================================
\section*{FORMACIÓN}

\entrada{Ingeniería Eléctrica -- Universidad Industrial de Santander}{En curso, 9.\textsuperscript{o} semestre}{%
Trabajo de grado: \href{https://www.zotero.org/groups/6637268/camacho-aranda-tesis-uis}{Estimación dinámica de la capacidad límite de integración de recursos energéticos distribuidos fotovoltaicos mediante PSO} (cotutela con Juan Camacho)}
{\small\textbf{Electivas en curso:} Mercados de energía eléctrica, Dimensionamiento fotovoltaico, Protecciones eléctricas. \textbf{Oyente:} Análisis de Sistemas de Distribución.}
\vspace{5pt}

\entrada{Técnico en Instalaciones Eléctricas Residenciales -- SENA}{Dic. 2021}{Cursado durante la educación media}

% ============================================================
% HABILIDADES
% ============================================================
\section*{HABILIDADES}
\fila{Potencia y energía}{Hosting Capacity, redes de distribución, sistemas fotovoltaicos, protecciones, mercados de energía, puestas a tierra}
\fila{Software y datos}{Python (NumPy, Seaborn), Rust (sistemas embebidos, en aprendizaje), OpenDSS, Scikit-learn, bases de datos relacionales, KiCad, Git, Linux y Bash, métodos numéricos}
\fila{Reportes técnicos}{Generación automatizada de reportes y presentaciones con Python y LaTeX, con gráficas de calidad profesional}
\fila{Automatización}{IEC 61131-3 (Ladder/FBD, GRAFCET), OpenPLC, Modbus TCP/IP}
\fila{Idiomas}{Español nativo \quad$\cdot$\quad Inglés en proceso de certificación (Altissia, UIS)}

% ============================================================
% PROYECTOS
% ============================================================
\section*{PROYECTOS}

\entrada{Sistema Automatizado de Clasificación y Empaque}{Curso de Accionamientos Eléctricos, UIS}{}
\vspace{-6pt}
\begin{itemize}
    \item Diseño eléctrico integral: cálculo de carga, motorreductores, protecciones y diagramas de potencia y control trifásico.
    \item Control en \textbf{IEC 61131-3} simulado en Factory IO vía Modbus TCP/IP (\href{https://youtube.com/playlist?list=PLUTAdDOdTlkNAtmTmna0_Z4TgAq-pZNoZ&si=ZV5tGop65upxRSAQ}{demostración en video}).
\end{itemize}
\vspace{4pt}

\entrada{BrotesPredict -- Alerta Temprana Epidemiológica}{Bootcamp de IA, proyecto final}{}
\vspace{-6pt}
\begin{itemize}
    \item Modelo \textbf{Gradient Boosting} que cruza calidad del agua (IDEAM) y morbilidad infantil (SIVIGILA) en 13 departamentos: \textbf{85.4\,\% de accuracy}, \textbf{0.912 de ROC-AUC}.
\end{itemize}
\vspace{4pt}

\entrada{PowerSense-3P -- Analizador de Redes Trifásico}{Proyecto personal, en desarrollo}{}
\vspace{-6pt}
\begin{itemize}
    \item Diseño de hardware con \textbf{ESP32 y ADC de 24 bits} y firmware en \textbf{Rust} para calidad de energía (RMS, FFT armónica, sags/swells; IEEE 519 e IEC 61000-4-30), con telemetría MQTT y clasificación de fallas.
\end{itemize}

% ============================================================
% EXPERIENCIA
% ============================================================
\section*{EXPERIENCIA}

\entrada{Cofundador y Gestor Operativo -- MentorUp, tutorías académicas}{2026}{}
\vspace{-6pt}
\begin{itemize}
    \item Coordiné junto a dos socios el enlace de estudiantes universitarios con tutores de su universidad y dicté directamente Métodos Numéricos, Sistemas de Control, Circuitos Eléctricos y Termodinámica.
    \item Promoví el servicio ante docentes y estudiantes de distintas áreas, y atendí clientes con necesidades muy diversas.
\end{itemize}
\vspace{4pt}

\entrada{Tutor Académico Independiente}{2024 -- Actualidad}{}
\vspace{-6pt}
\begin{itemize}
    \item Asesorías en asignaturas de ingeniería para estudiantes universitarios y de educación media.
\end{itemize}

% ============================================================
% CERTIFICACIONES
% ============================================================
\section*{CERTIFICACIONES}
\entrada{Bootcamp en Inteligencia Artificial, nivel básico (159 h)}{Dic. 2025}{MinTIC, Universidad de Antioquia, Universidad de Caldas y Ubicua Technology}
\entrada{Fundamentos de Programación, módulos I y II}{Dic. 2021}{MinTIC y Universidad Tecnológica de Pereira}

\end{document}
